\section{Kiểm Thử và Đánh Giá}

Như đã nêu ở phần Evaluation, nhóm xem việc đánh giá là một vòng phản
biện chạy song song với phát triển chứ không phải bước cuối. Chương này
trình bày cách kiểm thử thực tế trong giai đoạn pilot và đối chiếu kết
quả với các tiêu chí đã đặt ra.

\subsection{Phương Pháp Kiểm Thử}

Vì BeVietnam là một hệ thống nhiều thành phần (hai client, backend, AI
Core, vLLM), nhóm kiểm thử theo bốn lớp từ thấp đến cao, tóm tắt trong
Bảng~\ref{tab:test_levels}.

\begin{table}[H]
\centering
\caption{Bốn lớp kiểm thử của hệ thống}
\label{tab:test_levels}
\begin{tabularx}{\textwidth}{>{\raggedright\arraybackslash}p{3.0cm} X}
\toprule
\textbf{Lớp kiểm thử} & \textbf{Cách thực hiện} \\
\midrule
\textbf{Biên dịch / kiểu dữ liệu} & Kiểm tra web bằng trình biên dịch TypeScript, mobile bằng Gradle, mã Python của AI bằng \texttt{py\_compile}; đảm bảo toàn bộ build không lỗi trước khi tích hợp. \\
\textbf{Tích hợp dịch vụ} & Gọi trực tiếp các endpoint qua đường hầm public để xác nhận backend, proxy dịch vụ ngoài và AI Core hoạt động đúng từ ngoài mạng. \\
\textbf{Trải nghiệm trên thiết bị} & Cài bản APK thực lên điện thoại Android, kiểm tra bản đồ thời gian thực, định vị GPS và luồng chụp ảnh tại hiện trường. \\
\textbf{Hành vi AI} & Thử nghiệm Capture Judge với ảnh đúng và ảnh sai chủ thể để kiểm tra xem rubric có phân biệt được hay không. \\
\bottomrule
\end{tabularx}
\end{table}

\subsection{Kết Quả Kiểm Thử Chính}

Bảng~\ref{tab:test_results} liệt kê một số ca kiểm thử tiêu biểu và kết
quả quan sát được. Các ca tích hợp được thực hiện bằng cách gọi endpoint
qua đường hầm Cloudflare, mô phỏng đúng cách client thật truy cập hệ thống.

\begin{table}[H]
\centering
\caption{Một số ca kiểm thử tiêu biểu và kết quả}
\label{tab:test_results}
\begin{tabularx}{\textwidth}{>{\raggedright\arraybackslash}p{4.6cm} X >{\raggedright\arraybackslash}p{2.0cm}}
\toprule
\textbf{Ca kiểm thử} & \textbf{Kỳ vọng} & \textbf{Kết quả} \\
\midrule
POI thời gian thực tại trung tâm TP.HCM & Trả về danh sách quán/địa điểm quanh tọa độ, có phân loại nhóm. & Đạt \\
Kho nhiệm vụ \texttt{/question-pool} & Trả về toàn bộ nhiệm vụ cho Mission Path. & Đạt (133 nhiệm vụ) \\
Thời tiết khu vực & Trả về nhiệt độ, UV ước lượng, lượng mưa. & Đạt \\
Xác minh ảnh \textit{sai} chủ thể & Không được chấp nhận (rejected / needs\_review). & Đạt sau khi thay stub bằng mô hình thị giác \\
Build web (production) & Biên dịch và tạo trang \texttt{/storyline} thành công. & Đạt \\
Build mobile (APK debug) & Gradle hoàn tất, tạo được file cài đặt. & Đạt \\
\bottomrule
\end{tabularx}
\end{table}

Một quan sát quan trọng trong quá trình kiểm thử AI: phiên bản đầu của
Capture Judge chỉ kiểm tra ``có ảnh hay không'' rồi chấp nhận ở mức tin
cậy cao, nên một ảnh hoàn toàn không liên quan vẫn được duyệt. Đây chính
là loại lỗi mà chỉ kiểm thử hành vi mới phát hiện được, và là lý do nhóm
thay thế bằng pipeline rubric + ảnh tham chiếu + mô hình thị giác mô tả
ở chương trước.

\subsection{Đánh Giá theo Bốn Tiêu Chí}

Nhóm đối chiếu hệ thống đã hiện thực với bốn lớp tiêu chí đã đặt ra ở
chương System and Algorithm Design.

\textbf{Đúng chức năng.} Ba vòng lặp sản phẩm cốt lõi đều chạy thông từ
đầu đến cuối: người dùng mở bản đồ và thấy POI thời gian thực quanh mình;
mở Mission Path và đi qua các nút nhiệm vụ; chụp ảnh và nhận kết quả xác
minh có cơ sở. Các endpoint nền tảng đều phản hồi đúng khi gọi từ ngoài
mạng qua đường hầm.

\textbf{Chất lượng trải nghiệm.} Việc đồng bộ giao diện mobile theo hệ
màu Huế của web, bổ sung nhãn tên địa điểm trên bong bóng và nút định vị
nhanh giúp trải nghiệm liền mạch hơn giữa hai nền tảng. Mission Path dạng
lộ trình trực quan giúp người dùng nắm được tiến độ rõ ràng thay vì một
danh sách phẳng.

\textbf{Độ tin cậy của AI.} Sau khi thay stub bằng mô hình thị giác có
rubric, hệ thống không còn duyệt bừa ảnh sai chủ thể. Cơ chế
\textit{needs\_review} bảo đảm rằng khi mô hình không sẵn sàng hoặc còn
nghi ngờ, hệ thống không tự ý chấp nhận --- một hành vi an toàn quan trọng
trong pilot.

\textbf{Khả năng mở rộng.} Việc khóa nguồn POI về một endpoint dùng chung
cho cả hai client, tách nội dung nhiệm vụ (sinh offline) khỏi luồng điều
hướng, và đặt khóa dịch vụ ngoài sau một rate limiter chung khiến hệ
thống dễ nhân rộng sang địa phương khác mà không phải viết lại lõi.

\begin{note}[Giới hạn của kiểm thử]
Trong phạm vi pilot, kiểm thử còn nghiêng về kiểm tra tích hợp và hành vi
thủ công, chưa có bộ kiểm thử tự động đầy đủ hay đánh giá người dùng quy
mô lớn. Độ chính xác của mô hình thị giác cũng phụ thuộc vào chất lượng
ảnh tham chiếu lấy từ tìm kiếm web. Đây là những hướng cần củng cố khi
đưa hệ thống ra khỏi giai đoạn thí điểm.
\end{note}
